\section{Related Work}
\label{sec:related}

\katherine{Mention something about Genfer}
\katherine{Mention something about exact recursive probabilistic programming (another potential backend)}
\paragraph{Discrete-only probabilistic languages.}  
Several recent systems achieve \emph{exact inference} by restricting models to finite, discrete random variables. \emph{Roulette} extends Racket with first-class support for finitely-supported distributions and leverages symbolic evaluation to compile queries into weighted model-counting problems, enabling scalable exact conditioning on complex programs~\cite{Moy2025Roulette}. \emph{Dice} follows a similar philosophy in an OCaml DSL, compiling discrete programs to weighted model counting to perform exact Bayesian updates even on thousands of Boolean and finite-categorical variables~\cite{Holtzen2020Dice}. Earlier work in probabilistic logic programming, most prominently \emph{ProbLog}, annotates Prolog facts with probabilities and reduces inference to weighted Boolean formulas that can likewise be solved exactly~\cite{DeRaedt2007ProbLog}. These languages demonstrate the power of exact reasoning, but by construction they \emph{cannot express continuous random quantities}; consequently they offer no built-in path for analyzing models that naturally mix real-valued and discrete structure.  

\paragraph{Symbolic treatment of mixed models.}  
\emph{SPPL} occupies an intermediate point on the design spectrum. By enforcing syntactic restrictions that guarantee every program can be translated into a finite \emph{sum-product expression}, SPPL supports models with both discrete and continuous components while still admitting \emph{symbolic} exact inference~\cite{Saad2021SPPL}. Rather than discretizing the continuous parts, SPPL derives closed-form integrals when possible; exactness is retained for a narrow class of models. There is no mechanism for turning the remaining continuous structure into a discrete program that could reuse the powerful inference back-ends of Dice, Roulette, or ProbLog.

\paragraph{Continuous and hybrid PPLs with approximate inference.}  
The majority of widely-used PPLs treat \emph{continuous} distributions as first-class and rely on approximate inference. Systems such as \emph{Stan} (C++), \emph{PyMC} (Python), \emph{Pyro} (Python on PyTorch), \emph{TensorFlow Probability} and its precursor \emph{Edward} expose rich continuous distributions and obtain posteriors with Hamiltonian Monte Carlo, variational inference, or sequential Monte Carlo~\cite{Carpenter2017Stan,Salvatier2016PyMC3,Bingham2019Pyro,Dillon2017TFP,Tran2016Edward}. These frameworks deliver high accuracy for differentiable densities but can struggle with discrete latent structure or discontinuities that arise after naïve discretisation. Crucially, none of them offers an automated pathway to convert continuous sub-expressions into discrete replacements that would make exact inference feasible.

\paragraph{Universal languages supporting both regimes.}  
Universal or ``hybrid'' languages aim for expressiveness by embedding probabilistic primitives into general-purpose hosts. \emph{Anglican} (Clojure), \emph{WebPPL} (JavaScript), \emph{Figaro} (Scala) and \emph{Infer.NET} (C\#) permit arbitrary mixtures of discrete and continuous variables with inference based on importance sampling, Gibbs, expectation propagation, or lightweight MCMC~\cite{Tolpin2016Anglican,Goodman2014WebPPL,Pfeffer2009Figaro,Minka2018InferNET}. More recently, Julia-based systems such as \emph{Turing.jl} and \emph{Gen.jl} expose programmable inference interfaces, while \emph{Bean Machine} offers a declarative syntax atop PyTorch~\cite{Ge2018Turing,CusumanoTowner2019Gen,Tehrani2020BeanMachine}. Classic languages like \emph{Church} pioneered the ``evaluation-as-sampling'' view that underlies many of these systems~\cite{Goodman2008Church}. Discretization transformations are left entirely to the user and are not integrated with the inference engine.

\paragraph{Probabilistic program abstraction.}
The work of \cite{Holtzen2018Abstraction} is perhaps most similar to ours. They use abstract interpretation to transform hybrid programs into finite Boolean programs based on user-provided predicates, proving distributional soundness. However, their approach requires manual predicate selection and is limited to loop-free programs. In contrast, \Slice{} \emph{automatically infers cuts} from the program text using a type system and supports recursion and other advanced language features.

% \paragraph{Logical relations for probabilistic programs.}
% Our soundness proof builds on logical relations for probabilistic programs, including early work on denotational models~\cite{Kozen1981Semantics,Jones1989Probabilistic}, step-indexed relations for higher-order languages~\cite{Bizjak2015Step,Wand2018Contextual}, and program logics for verification~\cite{Barthe2009Formal,Barthe2012Privacy}. We adapt these techniques to prove discretization preserves semantics.

\paragraph{Positioning of our work.}  
Our contribution lies precisely at the intersection left open by the above lines of research. We propose a \emph{language-level transformation} that automatically discretizes continuous sub-programs while preserving the semantics required for exact Bayesian reasoning in the discrete fragment. When discretization successfully translates all continuous distributions, we can compile to any discrete PPL backend—we demonstrate with Dice and Roulette, but the approach generalizes to other discrete inference engines.

A key achievement is that when combined with Roulette as a backend, \Slice{} enables exact inference for continuous probabilistic programs with both higher-order functions and mutable state—a combination that no existing system can handle. While SPPL supports some continuous models exactly, it cannot handle higher-order functions or mutable state. While Roulette supports these features for discrete programs, it cannot handle continuous distributions. \Slice{} bridges this gap, discretizing continuous programs to leverage Roulette's advanced features.

Empirically, our approach achieves significantly faster runtimes compared to SPPL's symbolic methods. With Dice as a backend, we obtain maximal performance (up to 200× speedup) for first-order programs. With Roulette, we trade some performance for the ability to handle programs that no other exact inference system can solve. This flexibility—choosing between performance and expressiveness based on modeling needs—represents a new point in the design space of probabilistic programming systems.

\medskip  
In summary, existing discrete PPLs excel at exactness but lack continuous support, symbolic languages like SPPL require restrictive structure and are slower than model counting approaches, and continuous or hybrid frameworks favour approximate inference. To our knowledge none of these systems automates the conversion of continuous programs into a discrete representation amenable to exact inference; our work addresses precisely that missing piece.

\paragraph{Future work.}
Beyond our current results, several promising directions remain. First, we would like to extend our work to symbolically integrate tractable combinations of continuous distributions beyond the cumulative distribution function. Second, while we have preliminary evidence that partial discretization helps Monte Carlo methods converge significantly more quickly, a thorough empirical and theoretical investigation of this phenomenon would be valuable. Finally, we would like to improve mixed continuous-discrete inference by combining partial discretization with hybrid inference techniques, summing out discrete variables using weighted model counting while applying specialized continuous inference methods to the remaining continuous components.